% ============================================================================
%  data-fetcher-explained.tex
%  เอกสารอธิบาย src ของ Construction Data Fetcher แบบภาษาคนเขียน
%  Compile: xelatex (ต้องใช้ XeLaTeX เพราะมีฟอนต์ไทย)
% ============================================================================
\documentclass[12pt,a4paper]{article}

% ── ฟอนต์ไทย + อังกฤษ ผ่าน XeLaTeX ──────────────────────────────────
\usepackage{fontspec}
\usepackage{polyglossia}
\setdefaultlanguage{thai}
\setotherlanguage{english}

% ใช้ Sarabun (Google Fonts) จากโฟลเดอร์ fonts/ — ดีไซน์เดียวกับ TH Sarabun New
\setmainfont[
  Path=./fonts/,
  UprightFont=Sarabun-Regular.ttf,
  BoldFont=Sarabun-Bold.ttf,
  ItalicFont=Sarabun-Italic.ttf,
  BoldItalicFont=Sarabun-BoldItalic.ttf,
  Scale=1.25,
  Ligatures=TeX
]{Sarabun}

% mono font สำหรับโค้ด — ใช้ Menlo (มีใน macOS) + Thai fallback
\setmonofont[Scale=0.85]{Menlo}
\newfontfamily\thaifonttt[
  Path=./fonts/,
  UprightFont=Sarabun-Regular.ttf,
  BoldFont=Sarabun-Bold.ttf,
  Scale=1.15
]{Sarabun}

% ── การตัดคำไทย ─────────────────────────────────────────────────────
\XeTeXlinebreaklocale "th"
\XeTeXlinebreakskip = 0pt plus 1pt minus 0.1pt
\tolerance=2000
\emergencystretch=3em
\hyphenpenalty=5000

% ── หน้ากระดาษ ──────────────────────────────────────────────────────
\usepackage[a4paper,top=2.5cm,bottom=2.5cm,left=2.5cm,right=2.5cm]{geometry}
\usepackage{setspace}
\onehalfspacing

% ── สี + กล่อง + ตาราง + โค้ด ───────────────────────────────────────
\usepackage{xcolor}
\definecolor{ink}{HTML}{0A1628}
\definecolor{amber}{HTML}{D97706}
\definecolor{paper2}{HTML}{EDE5D3}
\definecolor{green}{HTML}{4D7C0F}
\definecolor{teal}{HTML}{0E7C7B}
\definecolor{rust}{HTML}{B94D2C}
\definecolor{codebg}{HTML}{F5F1E8}
\definecolor{linecol}{HTML}{2A4769}

\usepackage{tcolorbox}
\tcbuselibrary{skins,breakable}
\usepackage{booktabs}
\usepackage{tabularx}
\usepackage{array}
\usepackage{enumitem}
\usepackage{fancyhdr}
\usepackage{listings}
\usepackage{graphicx}

\usepackage{hyperref}
\hypersetup{
  colorlinks=true,
  linkcolor=ink,
  urlcolor=teal,
  pdftitle={Construction Data Fetcher — คู่มืออธิบาย src}
}

% ── ตั้งค่า listing สำหรับ Python ───────────────────────────────────
\lstdefinestyle{py}{
  language=Python,
  backgroundcolor=\color{codebg},
  basicstyle=\ttfamily\footnotesize,
  keywordstyle=\color{amber}\bfseries,
  commentstyle=\color{green}\itshape,
  stringstyle=\color{teal},
  numbers=left,
  numberstyle=\tiny\color{linecol},
  frame=single,
  rulecolor=\color{linecol},
  breaklines=true,
  showstringspaces=false,
  tabsize=2,
  captionpos=b,
  xleftmargin=2em,
  framexleftmargin=1.5em
}

% ── กล่องสไตล์ต่าง ๆ ────────────────────────────────────────────────
\newtcolorbox{infobox}[1]{
  colback=paper2, colframe=ink, boxrule=1pt,
  arc=2pt, title=#1, fonttitle=\bfseries,
  breakable, enhanced
}
\newtcolorbox{tipbox}{
  colback=codebg, colframe=amber, boxrule=1.5pt,
  arc=2pt, breakable, enhanced,
  borderline west={3pt}{0pt}{amber}
}

% ── header / footer ─────────────────────────────────────────────────
\pagestyle{fancy}
\fancyhf{}
\fancyhead[L]{\small\color{linecol}Construction Data Fetcher}
\fancyhead[R]{\small\color{linecol}คู่มืออธิบาย src}
\fancyfoot[C]{\thepage}
\renewcommand{\headrulewidth}{0.4pt}

% ── คำสั่งช่วย: หัวข้อแบบมีแถบ ──────────────────────────────────────
\newcommand{\barsection}[1]{%
  \section*{\color{ink}\rule[-0.2ex]{4pt}{2.2ex}\hspace{0.4em}#1}%
  \addcontentsline{toc}{section}{#1}%
}

\begin{document}

% ============================================================================
%  หน้าปก
% ============================================================================
\begin{titlepage}
\centering
\vspace*{2cm}
{\Huge\bfseries\color{ink} Construction Data Fetcher\par}
\vspace{0.5cm}
{\Large\color{amber} คู่มืออธิบายโค้ด (src) แบบเข้าใจง่าย\par}
\vspace{1.5cm}

\begin{tcolorbox}[colback=paper2,colframe=ink,width=0.85\textwidth,arc=4pt]
\centering\large
สคริปต์ Python ที่ดึง \textbf{ราคาวัสดุก่อสร้างจริง} \\
จาก API ของหน่วยงานภาครัฐไทย \\
แล้วบันทึกลงไฟล์ Excel อัตโนมัติ
\end{tcolorbox}

\vspace{2cm}
{\large
\begin{tabular}{rl}
\textbf{ไฟล์หลัก:} & \texttt{fetch\_construction\_data.py} \\[4pt]
\textbf{ภาษา:} & Python 3 \\[4pt]
\textbf{แหล่งข้อมูล:} & กระทรวงพาณิชย์ (MOC) + ฉลากเขียว \\[4pt]
\textbf{ผลลัพธ์:} & ไฟล์ Excel (.xlsx) หลาย sheet \\
\end{tabular}
\par}

\vfill
{\color{linecol} โครงการ Construction Cost Engine \\ ปรับปรุงล่าสุด: มิถุนายน 2569}
\end{titlepage}

\tableofcontents
\newpage

% ============================================================================
\barsection{1. ภาพรวม: สคริปต์นี้ทำอะไร}
% ============================================================================

ลองนึกภาพแบบนี้ครับ — เราอยากรู้ว่า ``ราคาวัสดุก่อสร้างในประเทศไทยตอนนี้เป็นยังไง''
ปกติเราต้องเข้าเว็บราชการหลายเว็บ คลิกหาเอกสาร PDF โหลดมาเปิดอ่านทีละไฟล์
ซึ่งเสียเวลามาก และข้อมูลก็กระจัดกระจาย

สคริปต์ตัวนี้ทำงานแทนเราทั้งหมด มันจะ:

\begin{enumerate}[leftmargin=2em]
  \item โทรไปถาม API ของหน่วยงานรัฐ 3 แห่ง (เหมือนโทรไปถามราคาทางโทรศัพท์ แต่เป็นคอมพิวเตอร์คุยกัน)
  \item รับข้อมูลกลับมาในรูปแบบ JSON (รูปแบบข้อมูลที่เครื่องอ่านง่าย)
  \item จัดระเบียบข้อมูลให้เป็นตาราง
  \item เซฟลงไฟล์ Excel ที่เราเปิดดูได้เลย พร้อมแยกเป็นหลาย sheet
\end{enumerate}

\begin{tipbox}
\textbf{จุดเด่นสำคัญ:} ข้อมูลทั้งหมดเป็น \textbf{ของจริง} ดึงสด ๆ จาก API ภาครัฐ
ไม่ใช่ข้อมูลปลอม (mock) ที่พิมพ์ขึ้นมาเอง — รันเมื่อไหร่ก็ได้ข้อมูลล่าสุดเมื่อนั้น
\end{tipbox}

\vspace{0.5em}
\begin{infobox}{ทำไมต้องเขียนสคริปต์ ไม่ทำมือ?}
เพราะราคาวัสดุ \textbf{เปลี่ยนทุกเดือน} ถ้าทำมือ เราต้องมานั่งโหลดใหม่ทุกครั้ง
แต่ถ้ามีสคริปต์ แค่รันคำสั่งเดียว \texttt{python3 fetch\_construction\_data.py}
ก็ได้ไฟล์ Excel ใหม่พร้อมข้อมูลล่าสุดทันที ภายในไม่กี่นาที
\end{infobox}

\newpage
% ============================================================================
\barsection{2. แหล่งข้อมูล: เราไปเอาข้อมูลจากไหน}
% ============================================================================

สคริปต์ดึงข้อมูลจาก API จริง 3 แหล่ง แต่ละแหล่งให้ข้อมูลคนละมุม:

\vspace{0.5em}
\begin{infobox}{แหล่งที่ 1 — CSI: ดัชนีราคาวัสดุระดับประเทศ}
\textbf{เจ้าของ:} กระทรวงพาณิชย์ (สำนักงานนโยบายและยุทธศาสตร์การค้า) \\
\textbf{URL:} \texttt{\small dataapi.moc.go.th/csi-indexes} \\
\textbf{ให้อะไร:} ตัวเลขดัชนีราคาวัสดุก่อสร้างทั้งประเทศ รายเดือน
ว่าเดือนนี้แพงขึ้นหรือถูกลงกี่เปอร์เซ็นต์เทียบปีก่อน \\
\textbf{เปรียบเทียบ:} เหมือน ``เทอร์โมมิเตอร์วัดราคา'' ของทั้งวงการก่อสร้าง
\end{infobox}

\begin{infobox}{แหล่งที่ 2 — CPIP: ดัชนีราคารายจังหวัด}
\textbf{เจ้าของ:} กระทรวงพาณิชย์ \\
\textbf{URL:} \texttt{\small dataapi.moc.go.th/cpip-indexes} \\
\textbf{ให้อะไร:} ดัชนีราคาก่อสร้างแยก \textbf{รายจังหวัด} (มากถึง 77 จังหวัด) รายเดือน
ใช้สำหรับงานก่อสร้างภาครัฐ \\
\textbf{เปรียบเทียบ:} เหมือนแผนที่ราคา บอกว่าจังหวัดไหนของแพง จังหวัดไหนของถูก
\end{infobox}

\begin{infobox}{แหล่งที่ 3 — Green Label: ราคาวัสดุจริงเป็นบาท}
\textbf{เจ้าของ:} ศูนย์ข้อมูลที่อยู่อาศัยแห่งชาติ (ฉลากเขียว) \\
\textbf{URL:} \texttt{\small nhicservices.nha.co.th/.../green\_label/mid} \\
\textbf{ให้อะไร:} ราคาวัสดุ \textbf{จริงเป็นบาท} กว่า 5,500 รายการ มียี่ห้อ สเปค หน่วย ครบ \\
\textbf{เปรียบเทียบ:} เหมือนใบเสนอราคาจริง ที่บอกว่า ``ปูนยี่ห้อนี้ถุงละเท่านี้''
\end{infobox}

\vspace{0.5em}
\begin{tipbox}
\textbf{ข้อควรรู้:} API ของกระทรวงพาณิชย์ (2 แหล่งแรก) \textbf{ตอบช้ามาก}
บางครั้งใช้เวลา 30--45 วินาทีต่อครั้ง สคริปต์เลยมีระบบ ``ลองใหม่'' (retry)
ถ้าครั้งแรกไม่สำเร็จ จะรออีกหน่อยแล้วลองอีก 2 ครั้ง
\end{tipbox}

\newpage
% ============================================================================
\barsection{3. โครงสร้างโค้ด: แต่ละส่วนทำหน้าที่อะไร}
% ============================================================================

โค้ดถูกแบ่งเป็นก้อน ๆ ให้อ่านง่าย แต่ละก้อนมีหน้าที่ชัดเจน
ลองดูภาพรวมก่อน แล้วค่อยเจาะทีละส่วน:

\vspace{0.5em}
\begin{center}
\begin{tabular}{>{\raggedright}p{5cm} p{9.5cm}}
\toprule
\textbf{ส่วนของโค้ด} & \textbf{หน้าที่ (พูดแบบบ้าน ๆ)} \\
\midrule
\texttt{class Config} & กล่องเก็บค่าตั้งต้น เช่น จะดึงปีไหน timeout กี่วินาที ปรับตรงนี้ที่เดียว \\
\addlinespace
\texttt{http\_get\_json()} & ``คนวิ่งของ'' โทรไปขอข้อมูลจาก API ถ้าพลาดก็ลองใหม่ให้ \\
\addlinespace
\texttt{class ConstructionDataFetcher} & ``หัวหน้างาน'' คุมการดึงทั้งหมด + เขียน Excel \\
\addlinespace
\quad\texttt{.fetch\_csi\_national()} & ดึงแหล่งที่ 1 (ดัชนีประเทศ) \\
\quad\texttt{.fetch\_cpip\_provincial()} & ดึงแหล่งที่ 2 (รายจังหวัด) \\
\quad\texttt{.fetch\_green\_label()} & ดึงแหล่งที่ 3 (ราคาจริง) \\
\quad\texttt{.write\_excel()} & รวมทุกอย่างเขียนลง Excel + จัดสวย \\
\addlinespace
\texttt{main()} & ``ปุ่มสตาร์ท'' สั่งให้ทุกอย่างทำงานตามลำดับ \\
\bottomrule
\end{tabular}
\end{center}

\vspace{1em}
เราจะไล่ดูทีละส่วนในหน้าถัด ๆ ไป พร้อมโค้ดจริงและคำอธิบาย

\newpage
% ============================================================================
\barsection{4. เจาะส่วนที่ 1: Config — กล่องตั้งค่า}
% ============================================================================

ส่วนนี้คือ ``แผงควบคุม'' ของสคริปต์ ถ้าอยากเปลี่ยนอะไร เปลี่ยนตรงนี้ที่เดียว
ไม่ต้องไปไล่แก้ทั้งโค้ด

\begin{lstlisting}[style=py,caption={คลาส Config — ค่าตั้งต้นทั้งหมด}]
class Config:
    OUTPUT_DIR = "output"      # โฟลเดอร์เก็บไฟล์ Excel
    TIMEOUT = 45               # รอ API นานสุดกี่วินาที
    RETRIES = 2                # ลองใหม่กี่ครั้งถ้าพลาด
    RETRY_DELAY = 3            # รอกี่วินาทีก่อนลองใหม่
    PACE_DELAY = 1.0           # หน่วงระหว่าง request กัน rate-limit

    CSI_FROM_YEAR = 2023       # ดึง CSI ตั้งแต่ปีไหน
    CSI_TO_YEAR = 2025         # ถึงปีไหน
    CPIP_FROM_YEAR = 2024
    CPIP_TO_YEAR = 2025
\end{lstlisting}

\begin{infobox}{อ่านง่าย ๆ}
อยากได้ข้อมูลย้อนหลังถึงปี 2020? เปลี่ยน \texttt{CSI\_FROM\_YEAR = 2020} จบ \\
API ตอบช้าจนพลาดบ่อย? เพิ่ม \texttt{TIMEOUT = 60} หรือ \texttt{RETRIES = 3} \\
ไม่ต้องเข้าใจโค้ดทั้งหมด แค่รู้ว่าแต่ละค่าหมายถึงอะไรก็ปรับได้
\end{infobox}

\newpage
% ============================================================================
\barsection{5. เจาะส่วนที่ 2: http\_get\_json — คนวิ่งของ}
% ============================================================================

ฟังก์ชันนี้มีหน้าที่เดียว: ``ไปขอข้อมูลจาก URL ที่กำหนด แล้วเอากลับมา''
แต่เพราะ API ราชการชอบตอบช้า/พลาด มันเลยมีระบบ \textbf{ลองใหม่อัตโนมัติ}

\begin{lstlisting}[style=py,caption={http\_get\_json — ขอข้อมูลพร้อม retry}]
def http_get_json(url, timeout, retries, retry_delay, logger=print):
    last_err = None
    for attempt in range(retries + 1):       # ลองทั้งหมด retries+1 ครั้ง
        try:
            req = urllib.request.Request(url, headers={...})
            with urllib.request.urlopen(req, timeout=timeout) as resp:
                raw = resp.read().decode("utf-8")
                return json.loads(raw)        # สำเร็จ! คืนข้อมูลเลย
        except (URLError, TimeoutError, JSONDecodeError) as err:
            last_err = err
            if attempt < retries:
                logger(f"retry {attempt+1}/{retries} ...")
                time.sleep(retry_delay)        # รอแล้วลองใหม่
    return None                                # ลองครบแล้วยังพลาด -> คืน None
\end{lstlisting}

\textbf{เล่าเป็นเรื่อง:} เหมือนเราโทรสั่งพิซซ่า ถ้าสายไม่ว่าง เราไม่ได้เลิกล้ม
เราวางสาย รอแป๊บ แล้วโทรใหม่ ลองสัก 3 ครั้ง ถ้ายังไม่ติดจริง ๆ ค่อยยอมแพ้
(คืนค่า \texttt{None} แปลว่า ``ไม่ได้ของ'')

\begin{tipbox}
\textbf{ทำไมต้องคืน None ไม่ใช่ crash?} เพราะถ้าแหล่งหนึ่งล่ม
เรายังอยากได้ข้อมูลจากอีก 2 แหล่งที่เหลือ การคืน \texttt{None} ทำให้สคริปต์
``ข้าม'' แหล่งที่พังไปได้ ไม่ทำให้ทั้งโปรแกรมตาย
\end{tipbox}

\newpage
% ============================================================================
\barsection{6. เจาะส่วนที่ 3: การดึงแต่ละแหล่ง}
% ============================================================================

หัวหน้างาน (\texttt{ConstructionDataFetcher}) มีลูกน้อง 3 คน แต่ละคนดึงคนละแหล่ง
หลักการเหมือนกันหมด: สร้าง URL $\rightarrow$ เรียก \texttt{http\_get\_json} $\rightarrow$
แปลงเป็นตาราง (DataFrame)

\begin{lstlisting}[style=py,caption={fetch\_csi\_national — ดึงดัชนีระดับประเทศ}]
def fetch_csi_national(self):
    url = (f"{self.MOC_BASE}/csi-indexes"
           f"?index_id=0000000000000000"      # 0000... = ดัชนีรวม
           f"&from_year={Config.CSI_FROM_YEAR}"
           f"&to_year={Config.CSI_TO_YEAR}")

    data = http_get_json(url, ...)             # เรียกคนวิ่งของ
    if not data:
        return None                            # ไม่ได้ก็ข้าม

    df = pd.DataFrame(data)                     # แปลง JSON -> ตาราง
    df = self._add_thai_date(df)               # เพิ่มคอลัมน์ period
    return df
\end{lstlisting}

\textbf{Green Label พิเศษหน่อย} เพราะข้อมูลที่ได้ซ่อนอยู่ลึกกว่า
ต้องเจาะเข้าไปที่ \texttt{result $\rightarrow$ items} และราคายังเป็นข้อความ
``2,731.10'' (มีลูกน้ำ) ต้องแปลงเป็นตัวเลขก่อนคำนวณได้:

\begin{lstlisting}[style=py,caption={แปลงราคาข้อความเป็นตัวเลข}]
# "2,731.10" (str) -> 2731.10 (float)
df["price_baht_numeric"] = (
    df["price_baht"].astype(str)
      .str.replace(",", "")             # เอาลูกน้ำออก
      .apply(lambda x: pd.to_numeric(x, errors="coerce"))
)
\end{lstlisting}

\newpage
% ============================================================================
\barsection{7. เจาะส่วนที่ 4: write\_excel — เขียนไฟล์ + จัดสวย}
% ============================================================================

พอดึงข้อมูลครบแล้ว ส่วนนี้รวมทุกอย่างเขียนลง Excel ไฟล์เดียว แยกเป็นหลาย sheet
แล้วจัด format ให้สวย (หัวตารางสีน้ำเงิน ตัวอักษรขาว ปรับความกว้างคอลัมน์อัตโนมัติ)

\begin{lstlisting}[style=py,caption={write\_excel — เขียนหลาย sheet}]
def write_excel(self, frames):
    ts = self.run_started.strftime("%Y%m%d_%H%M%S")
    filename = f"construction_data_{ts}.xlsx"   # ชื่อมี timestamp กันชนกัน

    with pd.ExcelWriter(filepath, engine="openpyxl") as writer:
        self._write_summary_sheet(writer, frames)   # sheet สรุป
        # เขียนแต่ละแหล่งเป็น sheet แยก
        csi.to_excel(writer, sheet_name="CSI_National")
        cpip.to_excel(writer, sheet_name="CPIP_Provincial")
        green.to_excel(writer, sheet_name="GreenLabel_Prices")
        # เขียน log การดึงทั้งหมด
        logs_df.to_excel(writer, sheet_name="Fetch_Logs")

    self._style_workbook(filepath)               # จัดสวยทีหลัง
\end{lstlisting}

\begin{infobox}{ทำไมชื่อไฟล์มี timestamp?}
เพราะเราอยากเก็บประวัติทุกครั้งที่ดึง ไม่ให้ทับกัน
ถ้ารันวันนี้ได้ \texttt{construction\_data\_20260610\_154817.xlsx}
รันพรุ่งนี้ก็ได้อีกไฟล์ชื่อใหม่ ทำให้ย้อนดูได้ว่าข้อมูลเดือนนี้ต่างจากเดือนก่อนยังไง
\end{infobox}

\newpage
% ============================================================================
\barsection{8. Data Pipeline: เส้นทางข้อมูลตั้งแต่ต้นจนจบ}
% ============================================================================

ภาพรวมว่าข้อมูล ``ไหล'' จากต้นทางถึงปลายทางอย่างไร:

\vspace{1em}
\begin{center}
\begin{tcolorbox}[colback=codebg,colframe=ink,width=\textwidth,arc=3pt]
\centering\ttfamily\small
\textbf{[1] API ภาครัฐ}\\
(CSI / CPIP / Green Label)\\
$\big\downarrow$ \textnormal{ส่ง request ผ่าน HTTPS}\\[4pt]
\textbf{[2] http\_get\_json()}\\
\textnormal{รับ JSON + retry ถ้าพลาด}\\
$\big\downarrow$ \textnormal{ได้ list/dict}\\[4pt]
\textbf{[3] pandas DataFrame}\\
\textnormal{แปลงเป็นตาราง + เพิ่มคอลัมน์ period + แปลงราคาเป็นตัวเลข}\\
$\big\downarrow$ \textnormal{ตารางสะอาด}\\[4pt]
\textbf{[4] write\_excel()}\\
\textnormal{รวมหลาย sheet + จัด format}\\
$\big\downarrow$ \textnormal{เขียนไฟล์}\\[4pt]
\textbf{[5] output/construction\_data\_<timestamp>.xlsx}\\
\textnormal{เปิดดูได้เลยใน Excel}
\end{tcolorbox}
\end{center}

\vspace{1em}
สรุปเป็นประโยคเดียว: \textbf{API $\rightarrow$ ดึง+ลองใหม่ $\rightarrow$ จัดตาราง
$\rightarrow$ เขียน Excel $\rightarrow$ เปิดดู}

\newpage
% ============================================================================
\barsection{9. ตัวอย่างจริง: Input / Process / Output}
% ============================================================================

มาดูตัวอย่างจริง ตั้งแต่สิ่งที่เราใส่เข้าไป จนถึงสิ่งที่ได้ออกมา

\subsection*{9.1 INPUT — สิ่งที่เราสั่ง}

\begin{lstlisting}[style=py,caption={แค่รันคำสั่งเดียว},numbers=none]
$ python3 fetch_construction_data.py
\end{lstlisting}

ไม่ต้องใส่ argument อะไรเลย ค่าทั้งหมดอยู่ใน \texttt{Config} แล้ว
(อยากเปลี่ยนปีก็ไปแก้ใน \texttt{Config} ก่อนรัน)

\subsection*{9.2 PROCESS — สิ่งที่เกิดขึ้นระหว่างรัน}

สคริปต์จะพ่น log ออกมาให้เห็นว่ากำลังทำอะไร (ตัวอย่างจริงจากการรัน):

\begin{lstlisting}[style=py,numbers=none,basicstyle=\ttfamily\scriptsize]
[1/3] CSI - ดัชนีราคาวัสดุก่อสร้างระดับประเทศ (รายเดือน)
      GET https://dataapi.moc.go.th/csi-indexes?index_id=00...
      retry 1/2 (TimeoutError) ...
      [OK] 24 แถว (2023/1 - 2024/12)
[2/3] CPIP - ดัชนีราคาก่อสร้างรายจังหวัด (77 จังหวัด)
      [OK] 728 แถว, 61 จังหวัด
[3/3] Green Label - ราคาวัสดุฉลากเขียวจริง (บาท)
      [OK] 5571 รายการ, 12 หมวดหมู่
\end{lstlisting}

\subsection*{9.3 OUTPUT — สิ่งที่ได้ (ข้อมูลจริง)}

\textbf{ตัวอย่างข้อมูล CSI ที่ได้} (ดัชนีราคาวัสดุระดับประเทศ):

\begin{center}
\begin{tabular}{cccc}
\toprule
\textbf{period} & \textbf{price\_index} & \textbf{yoy (\%)} & \textbf{ความหมาย} \\
\midrule
2023-01 & 113.1 & +3.5 & แพงขึ้น 3.5\% จากปีก่อน \\
2023-12 & 112.0 & $-$0.4 & ถูกลงเล็กน้อย \\
2024-12 & 112.4 & +0.4 & ทรงตัว \\
\bottomrule
\end{tabular}
\end{center}

\vspace{0.8em}
\textbf{ตัวอย่างข้อมูล CPIP ที่ได้} (ดัชนีรายจังหวัด, ม.ค. 2024):

\begin{center}
\begin{tabular}{ccc}
\toprule
\textbf{จังหวัด} & \textbf{period} & \textbf{price\_index} \\
\midrule
พังงา & 2024-01 & 106.7 \\
นนทบุรี & 2024-01 & 109.1 \\
เชียงราย & 2024-01 & 107.3 \\
\bottomrule
\end{tabular}
\end{center}

\vspace{0.8em}
\textbf{ตัวอย่างข้อมูล Green Label ที่ได้} (ราคาจริงเป็นบาท):

\begin{center}
\begin{tabular}{llcc}
\toprule
\textbf{ยี่ห้อ} & \textbf{วัสดุ} & \textbf{หน่วย} & \textbf{ราคา (บาท)} \\
\midrule
ซีแพค & วัสดุก่อ (คอนกรีต 400) & ลบ.ม. & 2,731.10 \\
ซีแพค & วัสดุก่อ (คอนกรีต 380) & ลบ.ม. & 2,667.40 \\
\bottomrule
\end{tabular}
\end{center}

\vspace{0.8em}
\textbf{ไฟล์ Excel ที่ได้} มี 6 sheet:

\begin{center}
\begin{tabular}{ll}
\toprule
\textbf{ชื่อ sheet} & \textbf{เนื้อหา} \\
\midrule
Summary & สรุปภาพรวมการดึง (เวลา, จำนวนแถวแต่ละแหล่ง) \\
CSI\_National & ดัชนีระดับประเทศ 24 แถว \\
CPIP\_Provincial & ดัชนีรายจังหวัด 728 แถว \\
GreenLabel\_Prices & ราคาวัสดุจริง 5,571 รายการ \\
GreenLabel\_Summary & สรุปราคา เฉลี่ย/ต่ำสุด/สูงสุด ต่อหมวด \\
Fetch\_Logs & log การดึงทุกขั้นตอน \\
\bottomrule
\end{tabular}
\end{center}

\newpage
% ============================================================================
\barsection{10. วิธีใช้งานจริง}
% ============================================================================

\subsection*{ขั้นที่ 1 — ติดตั้ง library ที่ต้องใช้}
\begin{lstlisting}[style=py,numbers=none]
pip install -r requirements.txt
# หรือ
pip install pandas openpyxl
\end{lstlisting}

\subsection*{ขั้นที่ 2 — รันสคริปต์}
\begin{lstlisting}[style=py,numbers=none]
cd data_fetcher
python3 fetch_construction_data.py
\end{lstlisting}

\subsection*{ขั้นที่ 3 — เปิดไฟล์}
ไฟล์จะอยู่ในโฟลเดอร์ \texttt{output/} ชื่อมี timestamp
เปิดด้วย Excel / Numbers / Google Sheets ได้เลย

\vspace{1em}
\begin{infobox}{สรุปสั้น ๆ}
สคริปต์นี้เปลี่ยนงานที่เคยต้องนั่งทำมือหลายชั่วโมง
(เข้าเว็บ โหลด PDF อ่าน จดใส่ Excel) ให้เหลือแค่ \textbf{รันคำสั่งเดียว}
แล้วได้ไฟล์ Excel พร้อมข้อมูลราคาวัสดุจริงจากภาครัฐ 3 แหล่ง
รวมกว่า 6,300 แถว ภายในไม่กี่นาที — และทำซ้ำได้ทุกเมื่อที่อยากได้ข้อมูลใหม่
\end{infobox}

\vfill
\begin{center}
\color{linecol}\rule{0.5\textwidth}{0.4pt}\\[0.5em]
\textit{จบคู่มือ — Construction Data Fetcher}
\end{center}

\end{document}
